\documentclass[11pt]{article}
\usepackage{amsmath,amssymb,amsthm,booktabs,array,longtable,geometry,hyperref,tikz}
\geometry{margin=1in}
\usetikzlibrary{arrows.meta,positioning,shapes}

\newtheorem{theorem}{Theorem}
\newtheorem{proposition}[theorem]{Proposition}
\newtheorem{lemma}[theorem]{Lemma}
\newtheorem{corollary}[theorem]{Corollary}
\newtheorem{definition}[theorem]{Definition}
\newtheorem{remark}[theorem]{Remark}

\title{CONJ\_BOUNDARY\_DECIDABLE: Exact Dependency Map\\
Between Schanuel, GAIL, EDB, and ELC Decidability}
\author{Monogate Research}
\date{2026-04-21}

\begin{document}
\maketitle

\begin{abstract}
CONJ\_BOUNDARY\_DECIDABLE asserts that membership in $\varepsilon(\mathbb{R})$
(the ELC field) is decidable, given sufficient hypotheses.
We produce the exact logical dependency map between Schanuel's conjecture (SC),
the Generalized Algebraic Isolation Lemma (GAIL), the Effective Depth Bound (EDB),
and decidability.
We catalog which classes of inputs already yield decidable membership
(via HLW, Baker's theorem, the Risch algorithm, and related tools)
and which remain open.
The main result: SC + EDB $\Rightarrow$ CBD (Conjectural Boundary Decidability),
and CBD $\Rightarrow$ a partial converse of SC.
EDB is identified as the critical missing bridge — the hypothesis with the
highest leverage for advancing decidability unconditionally.
\end{abstract}

\section{The Decision Problem}

\begin{definition}[ELC membership problem]
Given a real number $\alpha$ (presented as an oracle or by a finite algebraic/analytic description),
determine whether $\alpha \in \varepsilon(\mathbb{R})$ (the ELC field).
\end{definition}

\begin{remark}[Input model]
In practice, $\alpha$ is presented as:
(a) a finite algebraic expression over $\mathbb{Q}$ with $\exp, \ln, +, -, \times, \div$;
(b) an oracle that evaluates $\alpha$ to arbitrary precision; or
(c) a value defined by a convergent series or integral.
Case (a) is the most tractable; cases (b) and (c) are harder.
\end{remark}

\section{Key Hypotheses}

\begin{definition}[SC — Schanuel's Conjecture]
If $\alpha_1, \ldots, \alpha_n \in \mathbb{C}$ are $\mathbb{Q}$-linearly independent, then
\[
\mathrm{trdeg}(\mathbb{Q}(\alpha_1, \ldots, \alpha_n, e^{\alpha_1}, \ldots, e^{\alpha_n})/\mathbb{Q}) \geq n.
\]
\end{definition}

\begin{definition}[GAIL — Generalized Algebraic Isolation Lemma]
For all $\alpha \in \mathbb{R}$: $\alpha \notin \varepsilon(\mathbb{R}) \Rightarrow \sin(\alpha) \notin \varepsilon(\mathbb{R})$.
Equivalently: $\sin^{-1}(\varepsilon(\mathbb{R})) \cap \varepsilon(\mathbb{R})^c = \emptyset$.
\end{definition}

\begin{definition}[EDB — Effective Depth Bound]
For each $\alpha \in \varepsilon(\mathbb{R})$ with a computable presentation of complexity $C(\alpha)$,
there exists a computable function $D(C(\alpha))$ such that $\alpha$ is computable
by an F16 tree of depth at most $D(C(\alpha))$.
Equivalently: ELC membership at complexity $C$ is witnessed by a tree of bounded depth.
\end{definition}

\begin{definition}[CBD — Conjectural Boundary Decidability]
The ELC membership problem for inputs presented in complexity class $\mathcal{C}$
is decidable: there exists a Turing machine that, given input $\alpha \in \mathcal{C}$,
halts and outputs ``yes'' ($\alpha \in \varepsilon(\mathbb{R})$) or ``no'' ($\alpha \notin \varepsilon(\mathbb{R})$).
\end{definition}

\section{The Dependency Map}

\begin{theorem}[Main dependency chain]
\label{thm:chain}
The following implications hold (all arrows are strict — no known reversals):
\[
\mathrm{SC} \;\Rightarrow\; \mathrm{GAIL} \;\Rightarrow\;
\mathrm{CBD}(\text{trig args}) \;\xleftarrow{\text{also}}\; \mathrm{EDB}
\]
\[
\mathrm{SC} + \mathrm{EDB} \;\Rightarrow\; \mathrm{CBD}(\text{general})
\]
\[
\mathrm{CBD}(\text{general}) \;\Rightarrow\; \text{partial SC}
\]
\end{theorem}

\subsection{SC $\Rightarrow$ GAIL}

Proved in GAIL\_Investigation.tex (Theorem 4.1).
The key step: if $\alpha \notin \varepsilon(\mathbb{R})$ and $\sin(\alpha) \in \varepsilon(\mathbb{R})$,
then $e^{i\alpha} \in \varepsilon(\mathbb{C})$ (complex ELC field),
so $i\alpha \in \varepsilon(\mathbb{C})$ by taking log,
so $\alpha \in \varepsilon(\mathbb{R})$ — contradiction.
The step ``$e^{i\alpha} \in \varepsilon(\mathbb{C}) \Rightarrow i\alpha \in \varepsilon(\mathbb{C})$''
uses SC to prevent algebraic coincidences.

\subsection{GAIL $\Rightarrow$ CBD for trig arguments}

If GAIL holds, then for $\alpha \in \varepsilon(\mathbb{R})$, we know $\sin(\alpha) \in \varepsilon(\mathbb{R})$
(by direct computation). And for $\alpha \notin \varepsilon(\mathbb{R})$, GAIL gives
$\sin(\alpha) \notin \varepsilon(\mathbb{R})$. This makes the trig-evaluation step decidable.
Combined with HLW-based decidability for algebraic arguments, CBD holds for
``trig of ELC inputs.''

\subsection{SC + EDB $\Rightarrow$ CBD (general)}

Without EDB, even knowing SC, the decision procedure might not terminate:
there is no bound on how deep the ELC tree for $\alpha$ might be.
With EDB: given complexity $C(\alpha)$, search all F16 trees of depth $\leq D(C(\alpha))$.
If any evaluates to $\alpha$ (to sufficient precision), output ``yes''; otherwise ``no.''
SC ensures correctness of the classification; EDB ensures termination.

\subsection{CBD $\Rightarrow$ partial SC}

If CBD holds, it implies that questions of the form
``does $e^{i\alpha} \in \varepsilon(\mathbb{C})$ for $\alpha \in \varepsilon(\mathbb{R})$?''
are decidable. This gives a specific decidability result about transcendence degrees
of exponential extensions of $\mathbb{Q}$, which is a partial consequence of SC
in a restricted setting.

\section{Catalog of Decidable Cases}

\subsection{Always decidable (unconditional)}

\begin{center}
\renewcommand{\arraystretch}{1.2}
\begin{tabular}{lll}
\toprule
Input class & Decision & Reason \\
\midrule
$\alpha \in \mathbb{Q}$ & Yes (in $\varepsilon$) & $\mathbb{Q} \subset \varepsilon(\mathbb{R})$ by construction \\
$\alpha \in \overline{\mathbb{Q}}$ & Yes (in $\varepsilon$) & Algebraics computable by EPL \\
$\alpha = e^q$ for $q \in \mathbb{Q}$ & Yes (in $\varepsilon$) & $\exp(q)$ is 1 F16 node \\
$\alpha = \ln q$ for $q \in \mathbb{Q}_{>0}$ & Yes (in $\varepsilon$) & $\ln(q)$ is 1 F16 node \\
$\alpha = \sin(q)$ for $q \in \mathbb{Q} \setminus \{0\}$ & No (not in $\varepsilon$) & HLW: $e^{iq}$ transcendental \\
$\alpha = \cos(q)$ for $q \in \mathbb{Q} \setminus \{0\}$ & No (not in $\varepsilon$) & HLW \\
$\alpha = \tan(q)$ for rational $q/\pi \notin \mathbb{Q}$ & No (not in $\varepsilon$) & AIL theorem \\
$\alpha = \mathrm{erf}(q)$ for $q \in \mathbb{Q} \setminus \{0\}$ & No (not in $\varepsilon$) & T01 class (infinite zeros) \\
$\alpha = J_0(q)$ for $q \in \mathbb{Q} \setminus \{0\}$ & No (not in $\varepsilon$) & T01 class \\
$\alpha = \mathrm{Ai}(q)$ for $q \in \mathbb{Q}$ & No (not in $\varepsilon$) & T01 class \\
Linear ELC forms & Yes (in $\varepsilon$) & Baker's theorem (effective) \\
\bottomrule
\end{tabular}
\end{center}

\subsection{Decidable under Baker's Theorem}

Baker's theorem (1966, effective): for algebraic $\alpha_0, \alpha_1, \ldots, \alpha_n \neq 0$
and algebraic $\beta_j$ not all zero,
$|\beta_0 + \beta_1 \ln \alpha_1 + \cdots + \beta_n \ln \alpha_n| > C(\alpha,\beta) H^{-\kappa}$
where $H$ is the height and $C, \kappa$ are effective.

\begin{center}
\renewcommand{\arraystretch}{1.2}
\begin{tabular}{lll}
\toprule
Input type & Decision & Reason \\
\midrule
$\alpha = \beta_0 + \sum \beta_j \ln \alpha_j$ (algebraic $\alpha_j, \beta_j$) & Yes/No & Baker effective \\
$\alpha = $ period of an elliptic curve over $\overline{\mathbb{Q}}$ & Open & Conjectured not in $\varepsilon$; no proof \\
$\alpha = \Gamma(1/4)$ & Open & Nesterenko: algebraically independent from $\pi, e^\pi$; not ELC \\
$\alpha = \zeta(3)$ & Open & Apéry: irrational; ELC status unknown \\
\bottomrule
\end{tabular}
\end{center}

\subsection{Decidable via Richardson's Theorem (restricted)}

Richardson (1968) showed that the sign of expressions built from $\mathbb{Q}$,
$\pi$, $\ln 2$, $e^x$, $\sin(x \pi)$, and $|x|$ is undecidable in general.
However, for the restricted ELC field (no trig, no $\pi$), sign-decidability holds
for expressions with rational constants, by the Tarski--Seidenberg theorem applied
to the ordered field structure.

\begin{proposition}
For $\alpha$ presented as an F16 expression with rational constants
(i.e., $\alpha$ has an explicit ELC representation), ELC membership is trivially
decidable (``yes'' by the given representation).
The hard problem is the \emph{converse}: given a numeric value $\alpha$,
decide whether it \emph{has} an ELC representation.
\end{proposition}

\subsection{Open cases}

\begin{center}
\renewcommand{\arraystretch}{1.2}
\begin{tabular}{lll}
\toprule
Input & Decision status & Blocking hypothesis \\
\midrule
$\sin(\sin(1))$ & Open & WGAIL($\sin(1)$) = GAIL instance \\
$\sin^{(k)}(1)$ for $k \geq 2$ & Open & GAIL (k steps) \\
$e^{i\alpha}$ for $\alpha \in \varepsilon(\mathbb{R})$ & Open in general & SC \\
$\Gamma(1/3)$ & Open & Algebraic independence unknown \\
$\ln(\pi)$ & Open & Open unconditionally \\
$\pi + e$ & Open & Basic open problem \\
$\pi \cdot e$ & Open & SC would settle \\
Generic depth-$\geq 2$ ELC expression & Open & EDB \\
\bottomrule
\end{tabular}
\end{center}

\section{EDB: The Critical Missing Bridge}

\begin{proposition}[EDB leverage]
\label{prop:edb_leverage}
Of the three hypotheses SC, GAIL, EDB, the effective depth bound EDB has the
highest leverage for advancing decidability without full SC:
\begin{enumerate}
  \item \textbf{GAIL alone} $\Rightarrow$ CBD for trig-class inputs only.
    Trig is a specific sub-problem; most ELC membership questions do not involve sin/cos.
  \item \textbf{SC alone} $\Rightarrow$ GAIL (and hence the trig case),
    but SC does not give an algorithm without EDB (no termination bound).
  \item \textbf{EDB alone} $\Rightarrow$ CBD conditioned on correctness of the
    tree search (not on SC). If EDB holds and the search terminates,
    the decision is correct assuming ELC is what we think it is
    (no new algebraic coincidences). EDB plus precision bounds gives a
    practical decision procedure.
\end{enumerate}
\end{proposition}

\begin{definition}[EDB conjecture]
For all $\alpha \in \varepsilon(\mathbb{R})$ presented by a convergent limit
of rational F16 approximations with precision function $\varepsilon(n) = 2^{-n}$:
the minimum depth $D$ of an exact F16 tree for $\alpha$ satisfies
$D \leq f(K)$ where $K$ is the complexity of the presentation
and $f$ is a computable function.
\end{definition}

\begin{remark}[Why EDB is hard]
EDB would imply that there are no ``deeply buried'' ELC numbers —
no real number that is in $\varepsilon(\mathbb{R})$ but requires an
uncomputable depth of tree to witness.
Proving EDB would require understanding the full structure of the ELC field,
including all algebraic relations between $e$ and the reals.
It is at least as hard as Schanuel in the algebraic-independence direction.
\end{remark}

\section{Complete Dependency Diagram}

\[
\boxed{\mathrm{SC}}
\;\Longrightarrow\;
\boxed{\mathrm{GAIL}}
\;\Longrightarrow\;
\begin{cases}
\text{CBD for trig inputs} \\
\text{WGAIL}(s_1) \Rightarrow s_2 \notin \varepsilon(\mathbb{R}) \\
\text{Tower: } s_k \notin \varepsilon(\mathbb{R})\ \forall k \geq 1
\end{cases}
\]
\[
\boxed{\mathrm{EDB}}
\;\Longrightarrow\;
\text{CBD terminates (given correct classification)}
\]
\[
\boxed{\mathrm{SC} + \mathrm{EDB}}
\;\Longrightarrow\;
\boxed{\mathrm{CBD (general)}}
\;\Longrightarrow\;
\text{partial SC}
\]
\[
\boxed{\mathrm{HLW}}
\;\Longrightarrow\;
\text{CBD for algebraic-argument trig}
\quad (\text{proved; unconditional})
\]
\[
\boxed{\mathrm{Baker}}
\;\Longrightarrow\;
\text{CBD for log-linear forms over } \overline{\mathbb{Q}}
\quad (\text{proved; effective})
\]
\[
\boxed{\mathrm{Ax\text{-}Schanuel}}
\;\Longrightarrow\;
\text{function-field CBD (proved; does not specialize)}
\]

\section{Provable Decidability Results (Summary)}

\begin{theorem}[Decidable cases — unconditional]
The following instances of ELC membership are decidable without any open hypothesis:
\begin{enumerate}
  \item All algebraic numbers: in $\varepsilon(\mathbb{R})$ (constructively).
  \item All rational multiples of $e, e^2, \ln 2, \ln 3, \ldots, \ln p$ for prime $p$:
    in $\varepsilon(\mathbb{R})$ (constructively).
  \item $\sin(q)$, $\cos(q)$, $\tan(q)$ for nonzero rational $q$:
    \textbf{not} in $\varepsilon(\mathbb{R})$ (HLW).
  \item All special functions with infinitely many real zeros
    (erf, Bessel, Airy, Fresnel, lgamma on suitable domain):
    \textbf{not} in $\varepsilon(\mathbb{R})$ (T01 class theorem).
  \item $\beta_0 + \sum \beta_j \ln \alpha_j$ for algebraic $\alpha_j, \beta_j$:
    membership status decidable effectively (Baker).
\end{enumerate}
\end{theorem}

\begin{theorem}[Conditionally decidable cases]
Under GAIL (which follows from SC):
\begin{enumerate}
  \item $\sin^{(k)}(1) \notin \varepsilon(\mathbb{R})$ for all $k \geq 1$.
  \item More generally: $\sin(\alpha) \notin \varepsilon(\mathbb{R})$ for all $\alpha \notin \varepsilon(\mathbb{R})$.
  \item Same for $\cos, \tan$ (by analogous GAIL statements using $e^{i\alpha}$).
\end{enumerate}
\end{theorem}

\section{Research Priorities}

\begin{enumerate}
  \item \textbf{Unconditional WGAIL($s_1$)}: Decide $\sin(\sin(1)) \notin \varepsilon(\mathbb{R})$
    without SC. This is the single highest-priority open problem.

  \item \textbf{Formalize EDB for depth-1 expressions}:
    Prove that if $\alpha \in \varepsilon(\mathbb{R})$ and is ``naturally depth-1''
    (e.g., $e^q$ for $q \in \mathbb{Q}$), then no deeper expression equals $\alpha$
    unless algebraically forced. This is a weak form of EDB and may be accessible.

  \item \textbf{Richardson scope}: Precisely characterize which ELC sub-problems
    fall within Richardson's decidability zone (sign decision for sub-ELC expressions
    without $\pi$ or sin) versus which require Schanuel.

  \item \textbf{Four exponentials conjecture}: Determine whether the four-exponentials
    conjecture (weaker than SC, still open) suffices for GAIL in the specific
    case of $\alpha = \sin(1)$.
\end{enumerate}

\end{document}
